\documentclass[../main.tex]{subfiles}
\begin{document}
%==========================================TIÊU ĐỀ TẬP========================================
\begin{table}[h]
  \begin{adjustwidth}{-1cm}{}
    \begin{tabular}{>{\centering\arraybackslash}p{6cm}|>{\raggedright\arraybackslash}p{12.5cm}}
      \multirow{5}{6cm}
      % Cột bên trái: ÉPISODE và số tập
      {\\[-40pt]
        \flaregothic\fontsize{20pt}{20pt}\selectfont \centering
        {\contour{errorthree}{\textcolor{black}{ÉPISODE}}}\color{black}\\
        \burbank\fontsize{72pt}{72pt}\selectfont
        \includegraphics[height=3.12359cm]{../../../../Image/Book?/number/num11.png}
        \\[18pt]
        % \flaregothic\fontsize{14pt}{14pt}\selectfont
        % \color{epattr}BIENVENUE, \uline{\color{Banpire} Mel} !
        % \flaregothic\fontsize{18pt}{18pt}\selectfont
        % \color{epattr}ÉPISODE PERDU
        \color{black}
      }
       &
      % Dòng thứ nhất: tên tập tiếng Nhật
      {\fotskip\fontsize{18pt}{18pt}\selectfont \ruby{幽}{ゆう}\ruby{霊}{れい}\ruby{列}{れっ}\ruby{車}{しゃ}
      }                                             \\[6pt]
      % Dòng thứ hai: tên tập tiếng Anh
       & {\cafeteria\fontsize{18pt}{18pt}\selectfont Ghost Train}                                                                           \\[4pt]
      % Dòng thứ ba: tên tập tiếng Việt
       & {\vnmsans\fontsize{18pt}{18pt}\selectfont Con tàu ma}                                                                              \\[8pt]
      % Dòng thứ tư: Nhân vật xuất hiện
       & {\caviar{\fontsize{14pt}{14pt}\selectfont Track used: The Winter (La Serenissima - Igor Dvorkin, Duncan Pittock, \& Ellie Kidd.)}}
      \\[6pt]
      % Dòng cuối cùng: Thời gian diễn ra của tập (thường sẽ là ngày phát hành)
       & {\continuum\fontsize{16pt}{16pt}\selectfont Ngày phát hành: 29/04/2022}                                                            \\[8pt]
    \end{tabular}
  \end{adjustwidth}
\end{table}
%===========================================NỘI DUNG==========================================
\color{black}

\pov{Haragen Kyuen}

Mười giờ đêm.

Con đường dẫn tới ga Aogami về đêm vắng lặng đến mức tôi thấy lạnh cả sống lưng. Đèn đường thưa thớt, ánh sáng vàng mờ nhạt hắt xuống mặt đường như sắp tắt đến nơi. Tôi bước đi trước, Saikyou và Reiko sát bên, cả ba cùng nép vào nhau giữa màn đêm. Càng gần ga, tôi càng thấy rõ sự u ám kì lạ: tòa nhà gạch xám im lìm, chỉ có vài bóng đèn huỳnh quang yếu ớt lập lòe, như đang vật lộn chống lại bóng tối dày đặc.

Dù đồng hồ mới chỉ 10 giờ, nhưng không gian ở đây tĩnh lặng hệt nửa đêm. Không bóng người, không tiếng côn trùng, thậm chí cả tiếng gió cũng như bị nuốt chửng. Tôi dừng bước, hít một hơi, thấy hơi thở mình tỏa thành từng làn khói trắng. ``Sương đông'' - từ này chợt lóe lên trong đầu tôi, dẫu tôi chẳng chắc nó có thật hay không. Hơi lạnh luồn qua từng lớp áo, khiến tôi rùng mình.

Saikyou kéo sát cổ áo, khẽ càu nhàu bên tai tôi: ``Trời lạnh thật\dots\ mà anh có thấy không? Càng gần ga, càng lạnh, như kiểu chỗ này bị cách biệt với phần còn lại của thị trấn.''

``Như thể có một bức tường vô hình vậy,'' tôi thì thầm, cảm nhận rõ nhiệt độ giảm xuống theo từng bước chân. ``Không khí ở đây khác hẳn.''

Reiko ôm chặt khăn choàng mà tôi đưa cho lúc nãy, giọng run run: ``Em\dots\ em thấy gai người quá. Không giống mấy lần em đi học về muộn đâu. Ở đây\dots\ như không phải nơi dành cho người sống.''

``Em nói đúng,'' Saikyou khẽ gật đầu, mắt nhìn chăm chú vào những bóng đèn nhấp nháy. ``Có cảm giác như chúng ta đang bước vào một thế giới khác vậy.''

Tôi không đáp, chỉ siết nhẹ tay hai em như muốn nói ``Có anh đây''. Nhưng trong lòng tôi, nỗi bất an cũng đang lan ra. Những bóng đèn nhấp nháy phản chiếu trên tấm kính cửa ra vào, nơi tấm biển ``Ga Aogami (青上駅)'' đã bong tróc, càng khiến tôi thấy hoang tàn đến rợn người.

``Chúng ta nên quay lại không?'' Reiko hỏi nhỏ, giọng em như muốn biến mất vào không khí lạnh lẽo.

Tôi lắc đầu, cố giữ giọng vững vàng: ``Không, chúng ta đã đi xa thế này rồi. Phải tìm hiểu cho ra lẽ.''

Rồi, tôi ngẩng lên, và chợt nhìn thấy Akane-san.

\img{e2-3a}

Cô ấy đang bước vào trong ga, dáng đi chậm rãi, mái tóc sẫm màu khẽ lay theo gió đêm, cùng với cài tóc hình cục xương lệch sang trái. Cô vẫn ở trong bộ đồng phục. Khuôn mặt cô gần như không cảm xúc, đôi mắt trống rỗng, ánh sáng mờ nhạt khiến nét mặt càng u ám. Giữa bầu không khí lạnh lẽo này, Akane hiện ra như cái bóng mờ, vừa quen thuộc, vừa xa cách.

Tôi khựng lại, mắt mở to, tiếng gọi nghẹn trong cổ họng: ``\dots Akane-san? Cô ấy làm gì ở đây giờ này?''

Saikyou nhìn theo, giọng hạ thấp: ``Khoan\dots\ còn Yuka đâu? Hai người đó lúc nào cũng đi với nhau mà\dots\ Sao chỉ có một mình cậu ấy thôi vậy?''

``Kỳ lạ thật,'' tôi lẩm bẩm, cố gắng nhìn rõ hơn. ``Từ khi gặp họ, chưa bao giờ tôi thấy họ tách nhau ra cả.''

Reiko siết lấy tay áo tôi, run run: ``Onii-sama\dots\ có cảm thấy không? Trông cậu ấy\dots\ không giống Akane-san bình thường. Như thể cậu ấy\dots\ bị một thế lực nào đó hút vào vậy\dots''

``Như người mộng du,'' Saikyou thêm vào, giọng trầm xuống. ``Hoặc như bị thôi miên ấy.''

Tôi không trả lời. Trong khoảnh khắc, cả ba chúng tôi chỉ biết đứng lặng, nhìn bóng cô gái tóc nâu ấy dần khuất vào cửa ga, để lại sau lưng một khoảng không u tối và nỗi lạnh buốt dâng đầy trong lòng tôi.

``C-Còn chờ gì nữa, mau đuổi theo cô ấy!'' Saikyou hối thúc, kéo tay áo tôi. ``Nếu không nhanh lên, chúng ta sẽ mất dấu cô ấy mất!''

``Đúng vậy,'' tôi gật đầu, quyết tâm lấn át nỗi sợ. ``Dù có chuyện gì đi nữa, chúng ta cũng phải biết Akane-san đang làm gì.''

Chúng tôi lập tức lao tới cửa ga, vừa chạy vừa gọi to: ``''\textbf{Akane-san! Đợi đã!}''

Tiếng bước chân dồn dập vang lên trên nền gạch lạnh lẽo. Nhưng bất chợt, một tiếng thịch nhẹ, kèm với đó là tiếng ``Ái cha cha\dots''. Tôi giật mình khi nghe thấy tiếng kêu ngã, quay ngoắt lại cùng Saikyou, và thấy Reiko đang ngồi dưới đất, vừa va phải một cô gái.

\img{e2-3b}

``Đau, đau\dots\ Ờm\dots\ Cậu có sao không?'' cô em út lúng túng đỡ người kia.

Cô gái đó chống tay đứng dậy, chiếc máy ảnh vẫn cầm chắc trong tay, giọng trách móc: ``Đi đứng cẩn thận vào chứ. Cứ chạy như thế kẻo lại làm thương tích ai đó đấy.''

``Xin lỗi,'' tôi nói, bước tới giúp Reiko đứng dậy. ``Chúng tôi đang có chút vội.''

Reiko cúi đầu lia lịa, giọng nhỏ xíu: ``M-Mình xin lỗi\dots\ Tại chúng mình hơi vội\dots\ Cậu có đau ở đâu không?''

Saikyou và tôi bước nhanh tới, và ngay khoảnh khắc nhìn rõ mặt cô ấy, tôi thoáng sững người. Đôi tai thỏ trắng nổi bật trên mái tóc xanh nhạt, dài mềm, trong lọn tóc bím đuôi sam còn cài vài củ cà rốt nhỏ trông vừa buồn cười vừa lạ lùng. Ánh đèn huỳnh quang phản chiếu vào ống kính máy ảnh trên tay cô gái, sáng lấp lánh.

``Không thể nào\dots'' cô em lớn khẽ buột miệng.

``Cái gì vậy?'' tôi hỏi nhỏ, nghiêng người về phía Saikyou.

``Cô gái này\dots\ em có cảm giác đã gặp ở đâu đó rồi,'' em thì thầm, mắt không rời khỏi người lạ.

Reiko cũng nhìn chằm chằm, đôi mắt tròn xoe. Còn tôi, trong thoáng chốc, cảm giác rùng mình lan dọc sống lưng. Từ sáng tới giờ, những người chúng tôi gặp trong thị trấn đều mang vẻ gì đó khác lạ, còn cô gái này thì hoàn toàn ngược lại - khí chất cô toát ra giống hệt như\dots\ người bình thường.

Mà nhắc mới nhớ, đó có phải là cô gái vừa chạy ngang qua ba chúng tôi hồi nãy không nhỉ?

Cô gái tai thỏ nghiêng đầu, trông khá ngạc nhiên khi thấy cả ba chúng tôi cứ đứng nhìn mình như thể gặp ma.

``Ờm\dots\ có chuyện gì sao? Mặt mình có gì à?'' cô hỏi, giọng có chút ngập ngừng. ``Hay là\dots\ tai thỏ của mình trông kỳ cục lắm?'' Cô đưa tay lên chạm nhẹ vào đôi tai giả.

Tôi chưa kịp đáp thì hình ảnh cô nàng tai chó vụt sáng trong đầu. Chúng tôi cùng nhau giật mình nhớ ra mục tiêu ban nãy. Tôi liền quay phắt lại phía cửa ga, sốt ruột thốt lên một cách bộc phát: ``Chết tiệt, Akane-san!''

``Akane-san? Cậu cần gặp cậu ấy có việc gì sao?'' cô gái tai thỏ hỏi, vẻ mặt bỗng trở nên nghiêm túc.

Nhưng tôi không còn thời gian để trả lời. Cả ba chúng tôi định chạy tiếp, nhưng giọng cô gái tai thỏ vang lên phía sau, gọi giật lại: ``Khoan đã! Các cậu không cần phải vội như thế đâu. Chuyến tàu cuối cùng vừa mới rời đi rồi.''

Chúng tôi khựng lại trong một thoáng. Tôi ngoái đầu lại, cau mày: ``\dots Cái gì cơ?''

``Tàu cuối cùng?'' Saikyou hỏi lại, giọng đầy nghi hoặc. ``Cậu chắc chứ?''

Cô ấy gật nhẹ, giọng chắc nịch: ``Ở ga này, tàu kết thúc sớm lắm. Khoảng mười giờ là hết chuyến trong ngày. Giờ này chẳng còn tàu nào chạy đâu.''

``Nhưng chúng tôi vừa thấy Akane-san bước vào ga,'' tôi phản đối, cảm thấy có gì đó không ổn. ``Cô ấy không thể biến mất được.''

``Có lẽ cô ấy chỉ đi ngang qua thôi?'' cô gái tai thỏ đề xuất, nhưng giọng cô có chút lấp lửng.

Thông tin đó chỉ càng khiến bầu không khí lạnh lẽo hơn. Nhưng chưa ai trong chúng tôi kịp phản ứng, tôi đã cảm thấy một nỗi bất an dấy lên, thôi thúc mình không thể chần chừ thêm nữa.

``Dù vậy\dots'' tôi nói, nắm chặt tay, ``chúng tôi vẫn phải vào. Cô ấy ở trong đó.''

``Onii-sama có chắc không?'' Reiko hỏi, giọng run run. ``Em có linh cảm không hay về chuyện này.''

``Chúng ta đã đi xa thế này rồi,'' Saikyou nói, ánh mắt kiên định. ``Không thể quay đầu được nữa.''

Không một ai trong chúng tôi phản đối. Cả ba lại cùng nhau lao về phía cửa ga, bỏ mặc ánh mắt khó hiểu của cô gái tai thỏ đứng phía sau.

``Này, đợi đã!'' tôi nghe tiếng cô gái tai thỏ gọi với theo, nhưng chúng tôi không còn thời gian để chần chừ nữa.

\ct

Chúng tôi lao thẳng vào sảnh ga, tiếng giày dội lên nền gạch lạnh vang vọng, như thể đang chạy vào một nơi bỏ hoang từ lâu. Không gian trống trải đến rợn người, tiếng bước chân chúng tôi dội lại từ những bức tường ẩm thấp. Ánh đèn huỳnh quang nhấp nháy trên trần, tạo ra những vệt sáng chập chờn trên sàn gạch bám đầy bụi. Nhưng chỉ cần vài bước nữa, khung cảnh sân ga mở ra trước mắt.

Và\dots\ ngay khi đặt chân tới nơi, cả ba chúng tôi đã nhìn thấy nó.

\img{e2-2a1}

Tôi chết lặng.

Một con tàu hơi nước đang đứng chờ sẵn ở đó. Số hiệu 211493 lấp lóe ngay ở đầu máy, ánh đèn vàng nhạt từ các toa hắt ra, mờ đục nhưng đủ sáng để thấy rõ cửa đang mở rộng. Không khí nặng trĩu khói mỏng, mùi dầu máy len lỏi vào từng hơi thở. Thân tàu đen bóng phản chiếu ánh đèn ga yếu ớt, tạo ra những vệt sáng kỳ ảo như những linh hồn đang nhảy múa trên bề mặt kim loại.

``Không thể nào\dots''  tôi buột miệng, cảm thấy cổ họng khô khốc.

Câu nói vừa thoát ra, Saikyou đã run lên thấy rõ. Ánh mắt em trừng lớn, giọng lạc đi: ``Đúng nó rồi\dots\ Con tàu này\dots\ chính là nó\dots\ Con tàu trong giấc mơ đó! Và Yuka-san\dots em thấy cậu ấy bước lên tàu này!''

``Saikyou, khoan đã!'' Tôi vươn tay định giữ em lại, nhưng đã quá muộn.

Chưa kịp suy nghĩ, em đã lao vọt về phía toa tàu. Tôi nghe tiếng Reiko hốt hoảng gọi với: ``\textbf{Onee-sama! Đợi em với!!}''

``Reiko! Đừng--'' Tôi chưa kịp ngăn thì cô em út cũng đã phóng theo cô chị.

Chẳng còn cách nào khác, tôi đuổi theo, tim đập thình thịch trong lồng ngực. ``\textbf{Akane-san! Yuka-san!}'' tôi gọi lớn, giọng vang vọng khắp sân ga vắng lặng. Ba chúng tôi gần như cùng lúc đặt chân lên bậc cửa tàu.

*XOẠCH!*

Ngay giây đó, cánh cửa trượt khép rầm lại sau lưng, như nuốt chửng chúng tôi vào trong. Tôi quay phắt lại, đập mạnh vào cửa, nhưng nó đã khóa chặt như thể được hàn kín từ bên ngoài. Đoàn tàu rung lên, từng thanh sắt, từng mối nối kim loại rít chói tai, kéo dài thành âm thanh chẳng khác nào một tiếng thở dài đau đớn. Khói trắng phun ra từ ống khói đầu tàu, dày đặc đến mức tràn ngập cả sân ga, biến cảnh vật bên ngoài thành một màn mờ mịt.

Tiếng còi vang lên: một tiếng dài, chát chúa, rền rĩ như xé toạc không gian, âm hưởng ngân xa tựa như vọng lại từ một thời đại đã chết. Cảm giác quen thuộc đến đáng sợ, hệt như lời đồn chúng tôi đã nghe từ người trong thị trấn, hệt như những gì Saikyou từng kể lại về giấc mơ. Và giờ, cả ba chúng tôi đang ở trong nó, bằng xương bằng thịt.

``Đây không phải là một con tàu bình thường,'' Saikyou lẩm bẩm, mắt mở to nhìn xung quanh. ``Nó giống hệt trong giấc mơ của em\dots\ từng chi tiết một.''

Tôi quay sang nhìn nội thất bên trong. Thay vì cổ kính như dáng vẻ bên ngoài, khung cảnh lại hiện đại đến kỳ quặc. Hàng ghế bọc nệm xanh xếp ngay ngắn, mỗi cụm hai ghế đối diện nhau. Trên đầu là giá đựng hành lý, sáng loáng, y hệt như kiểu tàu mà chúng tôi đã ngồi ngày hôm qua khi tới thị trấn Aogami.

``Kỳ lạ quá,'' tôi lẩm bẩm, tay lướt nhẹ trên bề mặt ghế. ``Bên ngoài như tàu cổ từ thế kỷ trước, nhưng bên trong lại\dots''

``\dots như tàu hiện đại,'' Reiko tiếp lời, giọng em nhỏ xíu. ``Như thể nó tồn tại ở hai thời điểm cùng lúc vậy.''

\img{e2-3d}

Nhưng có một thứ chẳng ăn nhập chút nào: đèn báo điểm dừng kế tiếp trên đầu toa nhấp nháy loạn xạ, các ký tự méo mó, như thể hệ thống bị lỗi nghiêm trọng. Tôi nheo mắt cố đọc những chữ số và ký tự đang nhảy múa trên bảng điện tử: 「無限 - 終点 - 虚無 - 帰らない」 (Vô hạn - Điểm cuối - Hư vô - Không trở về).

``Các em nhìn kìa,'' tôi chỉ tay lên bảng hiển thị. ``Những chữ đó\dots\ không phải tên ga bình thường.''

Tôi nghe tiếng cô em út nuốt khan, bàn tay cô vô thức siết chặt vạt áo. ``Em không thích chỗ này,'' cô thì thầm. ``Cảm giác như chúng ta đang đi vào một nơi\dots\ không thuộc về thế giới này.''

Saikyou bước lại gần, ngón tay gõ nhẹ lên thành ghế như thể kiểm tra độ thật. ``Còn mùi nữa,'' em khẽ nhăn mũi. ``Mùi sắt và mùi nước mưa cũ, như thể toa tàu này vừa mới chạy qua một trận mưa axit.''

``Em ngửi thấy cả mùi khói bếp,'' Reiko thêm vào, giọng run run. ``Như kiểu có ai đó vừa nấu ăn xong, nhưng mùi lại lạnh ngắt.''

Tôi nhìn xuống sàn. Lớp cao su trải lối đi vẫn bóng loáng, nhưng dưới ánh đèn huỳnh quang, tôi thoáng thấy những vệt mờ hình xoáy trôn ốc, như bàn chân người kéo lê trong vũng bùn rồi bị lau vội. ``Các em để ý chưa?'' tôi khẽ nói. ``Không có bụi. Không một hạt bụi nào trên giá hành lý, không một vết mờ trên cửa sổ.''

``Như thể\dots'' Saikyou hạ giọng, ``ai đó vừa mới dọn dẹp để đón khách. Đón chúng ta.''

Reiko rùng mình, ôm hai vai. ``Em thấy như có ai đang nhìn chúng ta từ phía sau lưng vậy.''

Tôi quay ngoắt lại. Hành lang toa tàu trống trơn, chỉ có tiếng máy lạnh rít lên rồi tắt ngúm, để lại một khoảng lặng đặc quánh. Nhưng trong khoảng lặng ấy, tôi nghe thấy (hay tưởng rằng mình nghe thấy) một nhịp thở không thuộc về ba chúng tôi.

``Đừng đứng lại quá lâu,'' tôi thì thầm. ``Chúng ta cần tìm Akane-san trước khi\dots'' tôi không dám nói hết câu: trước khi con tàu quyết định chúng ta không được phép rời đi nữa.

Saikyou ngồi phịch xuống ghế gần đó, hai vai run rẩy nhưng ánh mắt lại sáng rực, như vừa chạm vào điều gì đó mà chị đã mòn mỏi tìm kiếm. ``Nhưng Yuka-san và Akane-san đang ở đâu đó trên con tàu này,'' em nói, giọng đầy quyết tâm. ``Chúng ta phải tìm họ.''

Đoàn tàu bắt đầu chuyển bánh, từng cú giật nhẹ dần biến thành chuyển động đều đặn. Qua khung cửa sổ, sân ga Aogami dần lùi xa, chìm vào màn sương mù dày đặc. Còn tôi\dots\ tim đập dồn dập, vừa hoang mang, vừa thấp thỏm, chẳng biết chuyến đi này sẽ đưa chúng tôi tới đâu.

``Dù có chuyện gì xảy ra,'' tôi nói, cố giữ giọng vững vàng nhất có thể, ``chúng ta sẽ cùng nhau vượt qua. Đừng rời xa nhau, được chứ?''

Hai em gái gật đầu, ánh mắt đầy lo lắng nhưng cũng không kém phần kiên định. Chúng tôi đã bước vào cuộc hành trình không lối thoát này, và giờ chỉ còn cách tiến về phía trước.

Chúng tôi vẫn còn chưa kịp trấn tĩnh sau khi bước lên con tàu quái lạ này, thì bất chợt từ đằng sau dãy ghế phía trước, một bóng người bật ra, giọng lanh lảnh.

``Cho mình xin kiểu ảnh nhé!''

*TÁCH!*

\img{e2-3e}

Ánh chớp lóe lên ngay trước mắt. Tôi theo phản xạ chồm người về phía trước, tim nhảy dựng trong lồng ngực. Reiko thì kêu khẽ một tiếng, còn Saikyou giật mình đến mức gần như bật dậy khỏi ghế.

Trước mặt chúng tôi, chẳng ai khác ngoài cô gái tai thỏ ban nãy ở bên ngoài ga. Chiếc máy ảnh vẫn còn đưa cao trước mặt, màn hình hiển thị khuôn mặt ba chúng tôi vừa bị chụp. Cô ta khẽ nheo mắt, miệng mỉm cười mãn nguyện khi ngắm nghía loạt ảnh.

Tôi lắp bắp, vừa ngờ vực vừa không biết phải hỏi gì trước: ``C-Cậu là\dots''

Cô gái ấy ngẩng lên, nụ cười hiền, có phần ngây ngô mà cũng rất tự tin: ``Gọi mình là Uzuki. Mizuta Uzuki.''

\chrint

\begin{wrapfigure}[14]{r}{6.5cm}
  \includegraphics[height=8cm]{../../../../Image/Book?/uzuki_human.png}
\end{wrapfigure}

Mizuta Uzuki (\ruby{水}{みず}\ruby{田}{た}\ruby{兎}{う}\ruby{月}{づき}, \ruby{Mi-dư}{Thủy}-\ruby{ta}{Điền}\ \ruby{U}{Thỏ}-\ruby{dư-ki}{Nguyệt}) - cô gái với niềm đam mê nhiếp ảnh mà chúng tôi gặp tại ga. Có lẽ điều khiến tôi chú ý nhất về cô ấy là vẻ bình thường đến kỳ lạ - người duy nhất trong thị trấn này không toát ra bầu không khí u ám hay kỳ quặc nào.

Thoạt nhìn, Uzuki-san có vẻ trầm tính, không dễ bộc lộ cảm xúc. Nhưng qua cách cô ấy tiếp cận chúng tôi, tôi nhận ra cô thực sự rất quan tâm đến mọi người. Theo lời cô ấy, cô thường lang thang khắp thị trấn Aogami cả ngày lẫn đêm với chiếc máy ảnh, chụp lại mọi khoảnh khắc thú vị. Cô không phân biệt đối tượng, sẵn sàng giao lưu với bất kỳ ai một cách bình đẳng.

Nhìn cách cô ấy thao tác với chiếc máy ảnh, tôi có thể đoán được Uzuki-san không phải tay mơ. Cô từng chia sẻ rằng những tác phẩm của mình đã giành chiến thắng trong nhiều cuộc thi, chứng tỏ kỹ năng và cảm thụ nhiếp ảnh của cô khá ấn tượng.

\chrint

``À, tiện thể thì, tai thỏ trên đầu mình chỉ là bờm cài thôi nhé. Còn mấy củ cà rốt trong tóc mà các cậu thấy\dots'' cô khẽ lắc hai bím tóc đuôi sam xanh nhạt, để lộ mấy cái bông nhồi hình cà rốt, ``\dots chỉ là đồ gắn chơi vì mình thích ăn cà rốt thôi.''

``V-Vậy à\dots'' tôi thở dài. Reiko tròn mắt. Saikyou thì cau mày, chưa hiểu nổi vì sao một người thế này lại có mặt ở đây. Tôi thì chỉ kịp thở ra, gắng kéo bản thân bình tĩnh.

Cô em lớn hỏi ngay, giọng vẫn còn cứng: ``Cậu\dots\ lên tàu bằng cách nào vậy? Cửa đã đóng lại rồi mà.''

Uzuki-san bật cười khúc khích, giơ ngón tay chỉ về phía sau: ``À, mình lên từ cửa sau thôi. Không khó lắm đâu.''

Tôi và hai đứa em liếc nhau, chẳng ai kịp phản ứng gì, thì cô ấy nghiêng đầu, tò mò: ``Mà ba cậu đang định tìm ai vậy?''

Hai cô em gái quay sang nhìn tôi, như chờ người mở lời. Tôi gật nhẹ, rồi giải thích: `Bọn tôi đang tìm Akane-san. Vừa nãy thấy cô ấy bước lên tàu này. Ngoài ra em gái lớn của tôi,'' tôi liếc sang Saikyou,``còn mơ thấy Yuka-san cũng ở trên con tàu này. Thế nên bọn tôi mới liều lên để tìm.''

Uzuki-san khẽ ``à'' một tiếng, nét mặt vừa như ngạc nhiên vừa như đã quen với chuyện kỳ quặc.

``Ra thế\dots''

Cô thu máy ảnh lại, rồi thong thả đi dọc hành lang toa tàu. Chúng tôi chẳng còn cách nào ngoài việc bước theo.

\ct

Trên đường, câu chuyện trở nên nhẹ nhàng hơn một chút. Uzuki-san giới thiệu mình thuộc câu lạc bộ Nhiếp ảnh, thường xuyên mang máy ảnh theo người để ghi lại bất kỳ khoảnh khắc nào bắt mắt.

\dots Và tréo ngoe thế nào, cô ấy cũng là thành viên duy nhất của câu lạc bộ đó. Tôi dám chắc bởi không phải ai cũng thạo sử dụng máy ảnh kỹ thuật số, kể cả từ hai cô em gái, những người thường xuyên chụp ảnh tự sướng.

Reiko nhìn chằm chằm vào chiếc máy ảnh trong tay cô ấy, khẽ nói: ``Vậy\dots\ cậu có vẻ thích chụp ảnh lắm nhỉ?''

Uzuki-san cười, gật gù: ``Ừm. Mình ở trong CLB Nhiếp ảnh mà. Hồi trước còn từng chụp cho Akane-san và Yuka-san nữa cơ.''

\img{e2-3f}

Chúng tôi khựng lại. Cô em út gần như ngay lập tức hỏi, giọng gấp gáp:
``Cậu\dots\ từng chụp cho họ? Vậy cậu biết Akane-san và Yuka-san sao?''

``Có chứ,'' cô ấy đáp tỉnh bơ, như thể đó là điều hiển nhiên. ``Hai người ấy là bạn cùng lớp với mình mà. Thỉnh thoảng chúng mình còn hay ngồi chung bàn ăn trưa nữa.''

Saikyou nín thở, hỏi chậm rãi: ``Vậy\dots\ cậu vẫn nhớ Yuka-san?''

Uzuki hơi nghiêng đầu, vẻ bối rối nhẹ lướt qua, nhưng rồi cô gật khẽ:
``Mình nhớ chứ. Dù cậu ấy hay ngủ trong lớp, tỏ vẻ làm biếng, nhưng mỗi lần đến lớp đều ngồi cạnh Akane-san, cậu ấy trông không khác nào một con người khác vậy.''

Tôi cảm thấy sống lưng mình lạnh buốt. Reiko há hốc miệng, ánh mắt rạng lên vừa mừng vừa hoang mang. Saikyou thì nghiêm nghị, bàn tay siết chặt mép ghế như muốn chắc chắn rằng đây không phải là ảo giác.

``Cậu\dots\ cậu có biết họ thân thiết đến mức nào không?'' tôi hỏi khẽ, cố giữ giọng không run.

Uzuki mỉm cười, ánh mắt xa xăm:
``Akane-san thì luôn lo lắng cho Yula-san, nhắc cậu ấy uống thuốc, đội mũ khi trời lạnh. Còn Yuka-san thì hay đưa Akane-san đi chụp ảnh, bảo rằng `cậu ấy đẹp nhất khi cười tự nhiên'. Mình từng chụp lại khoảnh khắc cậu ấy đang chỉnh lại tóc mai cho cô bạn ở hành lang trường\dots\ hai người nhìn nhau như thể cả thế giới chỉ còn lại đối phương vậy.''

Cô dừng lại một chút, rồi thêm nhỏ:
``Mình nghĩ\dots\ họ không chỉ là bạn thôi đâu.''

Không khí trong toa tàu như ngưng đọng. Shino-san đã từng nói: sẽ có những ngoại lệ. Và giờ, trước mắt chúng tôi, Mizuta Uzuki - cô gái tai thỏ kỳ lạ này - chính là một ngoại lệ, và cô ấy đang mang theo những mảnh ký ức mà cả thị trấn dường như đã đánh mất.

\ct

Tàu lắc nhẹ trên đường ray, tiếng kim loại vang lên từng nhịp đều đặn, nhưng trong toa lại tĩnh lặng đến kỳ lạ. Cảm giác như thời gian bị kéo dãn ra, mỗi giây trôi qua đều nặng nề và chậm chạp.

Tôi nheo mắt nhìn quanh, thấy Uzuki-san vẫn cầm chiếc máy ảnh treo trước ngực, ánh mắt cô quét qua khung cảnh như đang tìm kiếm điều gì đó. Trong đầu tôi lúc này chỉ hiện lên một gương mặt duy nhất - Akane-san. Hình ảnh cô bước vào ga tàu với gương mặt tối sầm vẫn ám ảnh tôi, đôi mắt nâu trống rỗng và dáng đi như bị thôi miên.

``Uzuki-san\dots'' tôi cất tiếng, cố gắng giữ giọng bình tĩnh nhưng đủ rõ để át tiếng bánh xe nghiến trên đường ray. ``Cậu có biết Akane-san đang ở đâu không?''

Cô ấy thoáng giật mình khi nghe tôi hỏi, rồi đưa ngón tay chỉ về phía cuối toa. Ánh mắt cô thoáng vẻ lo lắng.

``Mình nghĩ là\dots\ cậu ấy đi về hướng đó?''

Saikyou ngay lập tức cau mày, vẻ không tin tưởng hiện rõ trên gương mặt. ``Cậu biết chắc à?''

Uzuki-san chỉ nhún vai, lắc đầu với nụ cười gượng gạo: ``Mình cũng không chắc đâu\dots\ Nhưng mình có thấy bóng ai đó đi về hướng đó.''

Tôi liếc nhìn hai cô em gái, cả ba chúng tôi gần như cùng lúc trao đổi ánh mắt quyết tâm. Không cần nói thêm lời nào, chúng tôi đồng loạt rảo bước về phía cuối toa.

``Cảm ơn Uzuki-san nhé!'' tôi vội nói vọng lại, không quên quay đầu nhìn cô gái tai thỏ đang đứng sững, nghiêng đầu với vẻ khó hiểu.

``Đi đâu mà vội mà vàng vậy\dots'' tôi nghe tiếng cô ấy lẩm bẩm phía sau, chiếc máy ảnh rung nhẹ trong tay cô.

\img{e2-3h}

Tôi cũng muốn hỏi Uzuki-san kỹ hơn về mối quan hệ của họ, nhưng\dots\ có lẽ tôi sẽ để dịp khác vậy.

\ct

Chúng tôi bước qua cửa ngăn toa, tiếng kim loại kêu ken két khi tôi đẩy cánh cửa nặng nề. Và ngay lập tức, một khung cảnh khiến cả ba chúng tôi phải khựng lại. Phía cuối toa, Akane-san đứng quay lưng về phía chúng tôi. Bóng dáng cô im lìm, vai hơi run nhẹ như đang kìm nén một điều gì đó. Từ chỗ chúng tôi đứng, tôi chỉ có thể thấy mái tóc buông xuống và đôi tai chó rung khẽ theo từng nhịp thở.

Một luồng khí lạnh đột ngột xộc thẳng vào người tôi. Không phải gió từ cửa sổ, không phải hơi lạnh của đêm khuya. Một thứ sát khí, một cảm giác như có lưỡi dao vô hình treo lơ lửng ngay trên cổ. Tôi cảm nhận Reiko bên cạnh lặng lẽ siết chặt nắm tay, những đốt ngón tay trắng bệch. Saikyou, vốn mạnh miệng là thế, lần này chỉ nuốt khan, không thốt nên lời.

``\textit{Cái\dots\ cảm giác này\dots}'' tôi thầm nghĩ, từng thớ cơ trên người căng cứng như dây đàn. Bản năng mách bảo tôi nên quay đầu bỏ chạy, nhưng tôi không thể. Không phải lúc này.

Chúng tôi tới đây để tìm Akane-san, và đó là việc mà chúng tôi phải làm. Dù rằng có thể cô ấy đang tính bày trò gì đó mang ý hãm hại chúng tôi.

Tiến lại gần từng bước, tôi cảm nhận được từng bước chân nặng như bị đè đá. Tiếng giày chạm sàn tàu vang lên trầm đục, hòa với tiếng bánh xe lăn đều đặn. Không gian xung quanh như méo mó, kéo dài ra theo mỗi bước chân tôi tiến về phía trước. Cảm giác như đi trong một giấc mơ, nơi khoảng cách không còn tuân theo quy luật thông thường.

``Onii-sama\dots'' Reiko thì thầm sau lưng tôi, giọng em run rẩy. ``Em không thích điều này\dots\ Có gì đó không đúng với Akane-san.''

Saikyou nắm chặt tay em gái, giọng trầm xuống: ``Đừng sợ, Reiko. Chúng ta phải tìm hiểu chuyện gì đang xảy ra.''

Khi khoảng cách đã gần đủ để nghe rõ tiếng thở của nhau, tôi lấy hết can đảm, mở miệng: ``Akane-san\dots\ Yuka-san\dots\ cô ấy mất tích rồi!''

Không một lời báo trước, Akane-san lập tức xoay người. Đôi mắt nâu như cá ươn trợn lên, ánh nhìn sắc lạnh như mũi dao đâm thẳng vào tim tôi. Tôi cảm thấy cơ thể mình như bị đóng băng tại chỗ. Giọng nói của cô hạ xuống trầm khàn, rợn người, chỉ gói gọn đúng ba từ:

``Cậu đây rồi.''

Chừng đó cũng đủ khiến tôi cảm thấy như bị ai đó bóp nghẹt cổ họng. Khoảnh khắc đó, không gian bỗng trở nên nặng trĩu. Không khí đặc quánh, khiến từng hơi thở của tôi như nghẹn lại trong lồng ngực. Hơi lạnh hòa lẫn mùi tanh tưởng tượng, khiến cả ba chúng tôi đứng chôn chân tại chỗ, không kịp phản ứng.

``A-Akane-san, thực ra thì\dots'' tôi cố gắng nói, nhưng giọng tôi run rẩy đến mức khó nhận ra. Nhưng chưa để tôi nói câu gì tiếp theo\dots\

Đột ngột, ánh sáng trong toa méo mó. Cả không gian ngập chìm trong màu đỏ thẫm như máu, nhuộm từng hàng ghế, từng ô cửa sổ thành sắc kinh dị. Tiếng tàu nghiến rít càng vang dội, như tiếng gào khóc từ nơi xa xăm. Tôi cảm thấy Reiko nắm chặt lấy tay áo tôi, còn Saikyou thì lùi lại một bước, mắt mở to kinh hoàng.

``Cái gì vậy?'' Saikyou thốt lên. ``Chuyện quái gì đang xảy ra với con tàu này?''

Tôi cảm nhận được nhiệt độ trong toa đột ngột giảm xuống, hơi thở chúng tôi bỗng hóa thành những làn khói trắng mỏng. Tiếng rít của bánh xe tàu như tiếng gào thét của linh hồn đau đớn, xuyên thấu vào tận xương tủy.

\img{e2-3i}

``Akane-san\dots?'' một giọng nói run rẩy vang lên từ cửa. Uzuki-san bước vào, mắt mở to kinh hãi. Nhìn thấy người bạn tai chó vốn hiền hòa giờ đứng đó, toàn thân tỏa ra sát khí ngùn ngụt, khiến cô không kìm được mà toát mồ hôi lạnh. ``Cậu\dots\ cậu đang làm gì vậy? Chuyện gì đang xảy ra với cậu?''

``\textbf{Uzuki-san, đừng lại gần!}'' tôi hét lên, nhưng dường như đã quá muộn.

Akane-san quay phắt lại, đôi mắt như hai hố lửa, găm chặt vào tôi, Saikyou và Reiko. Cô tiến từng bước, giọng khàn đục vang vọng như đến từ nơi khác, một nơi không thuộc về thế giới này:

``Và lần này\dots\ cậu không thể thoát được đâu.''

Mỗi bước chân của Akane-san dường như làm rung chuyển cả toa tàu. Những bóng đen kỳ quái bắt đầu nhảy múa trên tường, méo mó và vặn vẹo như những linh hồn bị tra tấn. Reiko bật khóc, còn Saikyou đứng chắn trước em gái, cố tỏ ra can đảm dù toàn thân đang run rẩy.

``Akane-san, tỉnh lại đi!'' Uzuki-san gọi lớn, giọng đầy tuyệt vọng. ``Đây không phải là cậu! Cậu không phải là người như thế này!''

\end{document}